\chapter{L'apparato cardiovascolare}

% ---------------------------------------------------------------------------
\section{Generalità}

\textbf{Cuore}: muscolo, cavo, impari e paramediano; funzione di pompa, permette la circolazione
del sangue all'interno dei vasi sanguigni.

\rubrica{Funzioni dell'apparato cardiovascolare}
\begin{itemize}
  \item trasporto di sostanze nutritizie e di ossigeno, e loro distribuzione a tutte le parti
    dell'organismo;
  \item allontanamento dei prodotti del catabolismo cellulare ($CO_2$);
  \item termoregolazione;
  \item regolazione dell'omeostasi dei fluidi corporei;
  \item interviene nei processi immunitari, mediante il trasporto di cellule ad attività
    fagocitaria e di anticorpi.
\end{itemize}

% ---------------------------------------------------------------------------
\section{Topografia}

Il cuore è situato nella cavità toracica, nel mediastino medio, contenuto in un sacco
fibrosieroso detto \textbf{pericardio}. Ha forma di cono leggermente appiattito in senso
antero-posteriore, con la base rivolta in alto, a destra e all'indietro, e l'apice in basso, a
sinistra e in avanti.

\rubrica{Facce, base, apice, margini}
\begin{itemize}
  \item una faccia anteriore, detta \textbf{faccia sterno-costale} (per i suoi rapporti);
  \item una faccia postero-inferiore, detta \textbf{diaframmatica};
  \item una \textbf{base}, verso la quale convergono i grossi vasi del grande circolo e del
    piccolo circolo;
  \item un \textbf{apice} (o punta);
  \item 2 margini: uno destro (acuto), uno sinistro (ottuso).
\end{itemize}

% ---------------------------------------------------------------------------
\section{Configurazione esterna}

La superficie esterna è percorsa da solchi che segnano i limiti tra le cavità che costituiscono
il cuore.
\begin{itemize}
  \item \textbf{solco coronario} (atrioventricolare), a decorso trasversale: separa la porzione
    postero-superiore (atriale) da quella antero-inferiore (ventricolare);
  \item \textbf{solco interatriale}, a decorso longitudinale: esteso dal solco coronario alla
    cupola atriale, separa l'atrio destro e sinistro;
  \item \textbf{solchi interventricolari} anteriore e posteriore, a decorso longitudinale: estesi
    dal solco coronario fin quasi all'apice del cuore, segnano sulla faccia sterno-costale e sulla
    faccia diaframmatica il limite tra i due ventricoli;
  \item \textbf{solco terminale}: divide le 2 vene cave dall'auricola destra.
\end{itemize}

% ---------------------------------------------------------------------------
\section{Arterie coronarie}

\rubrica{Coronaria sinistra}
Il primo tratto è denominato \textbf{tronco comune}; dopo i primi 2--4 cm il vaso si divide in
due rami:
\begin{itemize}
  \item \textbf{arteria discendente anteriore}: decorre sulla parete anteriore del cuore;
  \item \textbf{arteria circonflessa}: decorre lungo la parete laterale del ventricolo sinistro.
\end{itemize}

\rubrica{Coronaria destra}
Irrora il ventricolo destro e porta il sangue alle pareti inferiore e posteriore del ventricolo
sinistro.

% ---------------------------------------------------------------------------
\section{Configurazione interna}

Il cuore ha \textbf{4 cavità}: 2 superiori, di forma cubica, denominate \textbf{atrio destro} e
\textbf{atrio sinistro}; 2 inferiori, di forma conica, dette \textbf{ventricolo destro} e
\textbf{ventricolo sinistro}.

La metà destra del cuore \textbf{non comunica} con la metà sinistra: sono separate da un setto,
per la maggior parte di natura muscolare, denominato superiormente \textbf{setto interatriale}
(divide i 2 atri) e inferiormente \textbf{setto interventricolare} (più spesso, separa i 2
ventricoli).
\begin{itemize}
  \item l'atrio destro comunica con il ventricolo destro attraverso l'\textbf{orifizio
    atrioventricolare destro}, provvisto di una valvola formata da 3 cuspidi: la \textbf{valvola
    tricuspide};
  \item l'atrio sinistro comunica con il ventricolo sinistro per mezzo dell'\textbf{orifizio
    atrioventricolare sinistro}, munito di una valvola costituita da 2 cuspidi affrontate: la
    \textbf{valvola mitrale} (o bicuspide).
\end{itemize}

\rubrica{Ventricolo sinistro}
Forma di cono leggermente appiattito.
\begin{itemize}
  \item parete anteromediale: corrisponde alla faccia sterno-costale e al margine sinistro del
    cuore;
  \item parete infero-mediale: corrisponde alla faccia diaframmatica e al setto
    interventricolare;
  \item base, apice, punta.
\end{itemize}

\rubrica{Atrio e ventricolo destro}
\begin{itemize}
  \item \textbf{cavità dell'atrio destro}: forma di un cubo; riceve il sangue refluo dal grande
    circolo tramite le 2 vene cave e il seno coronario;
  \item \textbf{cavità del ventricolo destro}: forma di una piramide triangolare, a base
    superiore e apice inferiore;
  \begin{itemize}
    \item parete anteriore: corrisponde alla maggior parte della faccia sterno-costale;
    \item parete postero-inferiore: corrisponde alla porzione destra, meno estesa, della faccia
      diaframmatica;
    \item parete mediale: formata dal setto interventricolare;
    \item apice, base.
  \end{itemize}
\end{itemize}

% ---------------------------------------------------------------------------
\section{Parete del cuore e pericardio}

\rubrica{Pericardio}
\begin{itemize}
  \item \textbf{pericardio fibroso}: strato più esterno;
  \item \textbf{pericardio sieroso}, in 2 foglietti: \textbf{pericardio parietale} e
    \textbf{pericardio viscerale} (o \textbf{epicardio}), tra i quali si trova la cavità
    pericardica.
\end{itemize}

\rubrica{Parete del cuore}
Dall'esterno all'interno: \textbf{epicardio}, \textbf{miocardio}, \textbf{endocardio}.

% ---------------------------------------------------------------------------
\section{Muscolatura cardiaca}

\begin{itemize}
  \item i \textbf{cardiomiociti} sono quasi totalmente dipendenti dalla respirazione aerobia per
    la sintesi di ATP necessaria alla contrazione; il sarcoplasma contiene centinaia di
    mitocondri e un'ingente quantità di \textbf{mioglobina}, per immagazzinare ossigeno;
  \item le membrane cellulari di due cardiomiociti sono unite da \textbf{desmosomi}, che
    contribuiscono a mantenere la posizione e la struttura tridimensionale del tessuto;
  \item a livello dei \textbf{dischi intercalari}, i cardiomiociti sono uniti da \textbf{giunzioni
    comunicanti}, che promuovono il passaggio di ioni e piccole molecole tra le cellule, creando
    una connessione elettrica diretta;
  \item i cardiomiociti sono connessi dal punto di vista meccanico, elettrico e chimico: il
    muscolo cardiaco è considerato un \textbf{sincizio funzionale}.
\end{itemize}

% ---------------------------------------------------------------------------
\section{Valvole cardiache}

\rubrica{Struttura generale}
\begin{enumerate}
  \item anello di tessuto connettivo, che entra nella costituzione dello scheletro fibroso del
    cuore;
  \item cuspidi di tessuto connettivo, che separano gli atri dai ventricoli;
  \item corde tendinee, che ancorano i margini delle cuspidi ai muscoli papillari della parete
    ventricolare.
\end{enumerate}

\rubrica{Lo scheletro fibroso del cuore}
\begin{itemize}
  \item costituito da 4 anelli di tessuto connettivo elastico, che circondano le 4 valvole
    cardiache;
  \item assicura sostegno fisico ai cardiomiociti, ai vasi e ai nervi;
  \item distribuisce le forze di contrazione;
  \item stabilizza la posizione delle valvole rispetto al miocardio ventricolare;
  \item isola fisicamente il miocardio atriale dal miocardio ventricolare.
\end{itemize}

\rubrica{Valvole semilunari}
\begin{itemize}
  \item costituite da tre spessi lembi semilunari connettivali, a concavità superiore;
  \item chiuse dalla pressione del sangue nelle arterie polmonari e nell'aorta quando i
    ventricoli sono rilasciati; si aprono per la pressione esercitata dal sangue quando i
    ventricoli si contraggono: la loro struttura previene il reflusso ventricolare del sangue.
\end{itemize}

\rubrica{Valvole atrioventricolari}
Costituite da una componente muscolare e da una fibrosa.
\begin{itemize}
  \item \textbf{valvola mitrale} (o bicuspidale): localizzata tra atrio e ventricolo sinistro,
    costituita da due lembi (cuspidi) e da un anello valvolare che congiunge le cuspidi alla
    parete del cuore;
  \item \textbf{valvola tricuspide}: localizzata tra atrio destro e ventricolo destro, forma
    ovale ($\phi$ 40 mm); ha la funzione di assicurare che il sangue non torni indietro nel suo
    percorso attraverso il cuore, ma che il flusso proceda in un'unica direzione, assicurando il
    corretto scambio di $O_2$ e $CO_2$ con i polmoni a vantaggio di tutto il corpo.
\end{itemize}

% ---------------------------------------------------------------------------
\section{Vascolarizzazione e innervazione del cuore}

\begin{itemize}
  \item \textbf{arterie}: coronarie;
  \item \textbf{vene}: il sangue refluo della circolazione coronarica viene raccolto da 3 sistemi
    venosi: \textbf{seno coronario} (grande e breve vaso, $\sim$3 cm), \textbf{vene cardiache
    anteriori}, \textbf{vene minime};
  \item \textbf{vasi linfatici};
  \item \textbf{nervi}: il cuore è innervato dal \textbf{plesso cardiaco}, con fibre
    parasimpatiche e simpatiche, sia afferenti sia efferenti; il plesso cardiaco è un complesso
    intreccio di rami nervosi e piccoli gangli parasimpatici, situato in corrispondenza della
    base del cuore.
\end{itemize}

% ---------------------------------------------------------------------------
\section{Ciclo cardiaco e sistema di conduzione}

\textbf{Ciclo cardiaco}: periodo compreso tra l'inizio di un battito cardiaco e l'inizio del
successivo; comprende periodi alternati di contrazione e rilasciamento, \textbf{sistole} e
\textbf{diastole}.

\rubrica{Sistema di conduzione}
\begin{itemize}
  \item \textbf{nodo senoatriale} (di Keith-Flack): origine dell'impulso;
  \item fasci internodali (anteriore, medio, posteriore): conducono l'impulso agli atri;
  \item \textbf{nodo atrioventricolare} (di Tawara-Aschoff);
  \item tronco comune del fascio atrioventricolare, che si divide in \textbf{branca destra} e
    \textbf{branca sinistra};
  \item \textbf{fibre di Purkinje}: distribuiscono l'impulso ai muscoli papillari e al miocardio
    ventricolare.
\end{itemize}

% ---------------------------------------------------------------------------
\section{Elettrocardiogramma}

La contrazione delle cellule muscolari cardiache è determinata dall'attivazione di un segnale
elettrico (potenziale d'azione), generato dall'attivazione dei canali ionici presenti sulla
membrana plasmatica. L'elettrocardiogramma (ECG) si registra sulla regione anteriore del torace,
sulla superficie cutanea, e rappresenta i cicli di depolarizzazione.

\rubrica{Onde e complessi}
\begin{itemize}
  \item onda \textbf{P}: depolarizzazione atriale;
  \item complesso \textbf{QRS}: depolarizzazione dei ventricoli;
  \item onda \textbf{T}: ripolarizzazione ventricolare.
\end{itemize}

\rubrica{Intervalli}
\begin{itemize}
  \item \textbf{P--Q}: tempo di conduzione dell'impulso dagli atri ai ventricoli;
  \item \textbf{S--T}: fine della depolarizzazione e inizio della ripolarizzazione ventricolare;
  \item \textbf{Q--T}: intero ciclo di depolarizzazione/ripolarizzazione ventricolare.
\end{itemize}

\rubrica{Principali alterazioni}
Fibrillazione, blocchi, tachicardia, bradicardia.

% ---------------------------------------------------------------------------
\section{Pressione arteriosa}

I valori della pressione si ottengono misurando la forza che il sangue esercita contro le pareti
delle arterie, in due diverse fasi del ciclo cardiaco. Si utilizza un apparecchio dotato di un
manicotto gonfiabile per il braccio e un misuratore di pressione (\textbf{sfigmomanometro}); si
ottengono così due valori, espressi in millimetri di mercurio (mmHg).

\rubrica{Misurazione (metodo auscultatorio)}
\begin{enumerate}
  \item il manicotto viene posto attorno alla parte superiore del braccio e gonfiato in modo da
    comprimere l'arteria brachiale e bloccare il flusso sanguigno: uno stetoscopio posto
    sull'arteria brachiale, distalmente al manicotto, non percepisce alcun rumore;
  \item la pressione nel manicotto viene lentamente diminuita, fino a quando, attraverso lo
    stetoscopio, si avverte un rumore detto \textbf{tono di Korotkoff} (creato dal flusso
    turbolento nell'arteria compressa); la pressione a cui si avverte il primo tono rappresenta la
    \textbf{pressione sistolica} (la più alta presente nell'arteria);
  \item la pressione a cui i toni scompaiono rappresenta la \textbf{pressione diastolica}: le
    arterie non sono più compresse e il flusso sanguigno, non più turbolento, è silenzioso.
\end{enumerate}

% ---------------------------------------------------------------------------
\section{Conseguenze del mancato apporto di sangue}

\begin{itemize}
  \item necrosi cellulare;
  \item placche arteriosclerotiche;
  \item shock;
  \item edema polmonare;
  \item aritmie;
  \item ischemie di altri organi.
\end{itemize}

\rubrica{Fattori di rischio}
Età, sesso, alimentazione, familiarità, stili di vita, ambiente, diabete, droghe, alcol.
